Sea la sucesión de Fibonacci definida por
\[
F(1)=1,\qquad F(2)=2,
\]
y recordemos que para cualquier \(n \ge 2\),
\[
F(n+1)=F(n)+F(n-1).
\]
La idea es representar dos términos consecutivos de la sucesión como un vector.
\[
\begin{bmatrix}
F(n) & F(n-1)
\end{bmatrix}
\]
Ahora nos hacemos la siguiente pregunta:
\textit{¿Existe una matriz que transforme este vector en el siguiente par de términos?}
Es decir, queremos encontrar una matriz \(M\) tal que
\[
\begin{bmatrix}
F(n) & F(n-1)
\end{bmatrix}
M
=
\begin{bmatrix}
F(n+1) & F(n)
\end{bmatrix}.
\]
Sabemos que
\[
F(n+1)=F(n)+F(n-1),
\]
por lo que el primer elemento del nuevo vector debe ser la suma de ambos elementos, mientras que el segundo simplemente debe conservar el valor de \(F(n)\).
La matriz que realiza exactamente esta transformación es
\[
M=
\begin{bmatrix}
1 & 1\\
1 & 0
\end{bmatrix}.
\]
En efecto,
\[
\begin{aligned}
\begin{bmatrix}
F(n) & F(n-1)
\end{bmatrix}
\begin{bmatrix}
1 & 1\\
1 & 0
\end{bmatrix}
&=
\begin{bmatrix}
F(n)\cdot1+F(n-1)\cdot1 &
F(n)\cdot1+F(n-1)\cdot0
\end{bmatrix}
\\[2mm]
&=
\begin{bmatrix}
F(n)+F(n-1) &
F(n)
\end{bmatrix}
\\[2mm]
&=
\begin{bmatrix}
F(n+1) &
F(n)
\end{bmatrix}.
\end{aligned}
\]

Es decir, cada multiplicación por la matriz avanza exactamente un término en la sucesión.

Por ejemplo,

\[
\begin{bmatrix}
2 & 1
\end{bmatrix}
\begin{bmatrix}
1 & 1\\
1 & 0
\end{bmatrix}
=
\begin{bmatrix}
3 & 2
\end{bmatrix},
\]

que corresponde a

\[
(F(2),F(1))
\longrightarrow
(F(3),F(2)).
\]

Multiplicando nuevamente,

\[
\begin{bmatrix}
3 & 2
\end{bmatrix}
\begin{bmatrix}
1 & 1\\
1 & 0
\end{bmatrix}
=
\begin{bmatrix}
5 & 3
\end{bmatrix},
\]

obtenemos

\[
(F(3),F(2))
\longrightarrow
(F(4),F(3)).
\]

Cada multiplicación avanza un paso más:

\[
(F(2),F(1))
\rightarrow
(F(3),F(2))
\rightarrow
(F(4),F(3))
\rightarrow
(F(5),F(4))
\rightarrow
\cdots
\]

Después de realizar \(n-2\) transformaciones, llegamos a

\[
\boxed{
\begin{bmatrix}
F(n) & F(n-1)
\end{bmatrix}
=
\begin{bmatrix}
F(2) & F(1)
\end{bmatrix}
M^{\,n-2}
}
\]

o, sustituyendo los valores iniciales,

\[
\boxed{
\begin{bmatrix}
F(n) & F(n-1)
\end{bmatrix}
=
\begin{bmatrix}
2 & 1
\end{bmatrix}
\begin{bmatrix}
1 & 1\\
1 & 0
\end{bmatrix}^{n-2}
}
\]

Por ejemplo, para calcular \(F(4561)\) solo necesitamos obtener

\[
M^{4559}.
\]

En lugar de multiplicar la matriz \(4559\) veces, utilizamos **Exponenciación Binaria**, que calcula la potencia de una matriz en únicamente \(O(\log n)\) multiplicaciones.

Como el tamaño de la matriz es constante (\(2\times2\)), cada multiplicación cuesta \(O(1)\). Por tanto, el algoritmo completo tiene complejidad

\[
\boxed{O(\log n)}.
\]